% ===== PRÉAMBULE CANONIQUE — à coller VERBATIM en tête de CHAQUE .tex =====
\documentclass[11pt,a4paper]{article}
\usepackage[utf8]{inputenc}\usepackage[T1]{fontenc}\usepackage{lmodern}
\usepackage[french]{babel}
\usepackage{amsmath,amssymb,mathtools}\usepackage{icomma}
\usepackage{siunitx}\sisetup{locale=FR}  % \qty{}{}, \si{}, \ang{}
\usepackage{xcolor}\usepackage{colortbl}
\usepackage{tikz}\usepackage{pgfplots}\pgfplotsset{compat=1.18}
\usepackage[siunitx,europeanresistors]{circuitikz} % circuits (élec.)
\usepackage[most]{tcolorbox}\usepackage{enumitem}\usepackage{array,booktabs}
\usepackage[a4paper,margin=2.2cm]{geometry}
\usetikzlibrary{arrows.meta,calc,angles,quotes,positioning,patterns,%
  backgrounds,shapes.geometric,decorations.pathmorphing,decorations.markings,intersections}
\definecolor{primary}{RGB}{222,49,99}\definecolor{primarydark}{RGB}{150,22,55}
\definecolor{defcol}{RGB}{0,114,178}\definecolor{methcol}{RGB}{52,92,148}
\definecolor{excol}{RGB}{0,135,90}\definecolor{attncol}{RGB}{200,40,55}
\tcbset{boxbase/.style={enhanced,breakable,boxrule=0.9pt,arc=2.5pt,
  left=3mm,right=3mm,top=2.3mm,bottom=2.3mm,fonttitle=\bfseries,coltitle=white}}
% --- boîtes TUTORIEL (noms EXACTS, ne pas renommer) ---
\newtcolorbox{cle}[1][]{boxbase,colback=defcol!6,colframe=defcol,
  title={\textsc{À retenir}\ifx\relax#1\relax\relax\else\textmd{ : \itshape #1}\fi}}
\newtcolorbox{methode}[1][]{boxbase,colback=methcol!7,colframe=methcol,sharp corners,
  title={\textsc{Méthode}\ifx\relax#1\relax\relax\else\textmd{ --- \itshape #1}\fi}}
\newtcolorbox{exemplebox}[1][]{boxbase,colback=excol!5,colframe=excol,
  title={\textsc{Exemple}\ifx\relax#1\relax\relax\else\textmd{ : \itshape #1}\fi}}
\newtcolorbox{piege}[1][]{boxbase,colback=attncol!6,colframe=attncol,
  title={\textsc{Piège}\ifx\relax#1\relax\relax\else\textmd{ : \itshape #1}\fi}}
\newtcolorbox{remarque}[1][]{boxbase,colback=black!4,colframe=black!45,
  title={\textsc{Remarque}\ifx\relax#1\relax\relax\else\textmd{ : \itshape #1}\fi}}
% --- boîtes EXERCICE / CORRIGÉ (noms EXACTS, auto-comptées) ---
\newtcolorbox[auto counter]{exo}[1][]{boxbase,colback=primary!4,colframe=primary,
  title={\textsc{Exercice}~\thetcbcounter\ifx\relax#1\relax\relax\else\textmd{ --- \itshape #1}\fi}}
\newtcolorbox[auto counter]{corrige}[1][]{boxbase,colback=excol!4,colframe=excol,
  title={\textsc{Corrigé}~\thetcbcounter\ifx\relax#1\relax\relax\else\textmd{ --- \itshape #1}\fi}}
\newcommand{\vv}[1]{\vec{#1}}\newcommand{\ds}{\displaystyle}
\newcommand{\R}{\mathbb{R}}\newcommand{\N}{\mathbb{N}}\newcommand{\Z}{\mathbb{Z}}
% =========================================================================
\begin{document}
\sisetup{per-mode=symbol}

\begin{corrige}[MCU : lesquelles sont fausses ?]
Fausses : \textbf{B et E}.
\begin{itemize}[leftmargin=1.4em,topsep=2pt,itemsep=1.5pt]
  \item \textbf{B (fausse)} : $\vv F_c$ pointe vers le centre, $\vv v$ est tangente : elles sont \emph{perpendiculaires}, pas de même sens.
  \item \textbf{E (fausse)} : $F_c=\tfrac{mv^2}{R}\propto v^2$ ; pour doubler $v$, il faut \emph{quadrupler} $F_c$.
  \item Correctes : A ($a_t=0$ en MCU) ; C ($\vv v\perp\vv F_c$) ; D ($F_c=m\omega^2R$ dépend de $m$, $\omega$, $R$) ;
        F (fil cassé $\Rightarrow$ trajectoire rectiligne tangente, par inertie).
\end{itemize}
\end{corrige}

\begin{corrige}[rotation d'un disque]
L'affirmation incorrecte est \textbf{B}.
\begin{itemize}[leftmargin=1.4em,topsep=2pt,itemsep=1.5pt]
  \item \textbf{B (fausse)} : la vitesse angulaire $\omega$ est la \emph{même} pour tous les points du disque
        (y compris au centre) ; elle n'est pas nulle en $O$.
  \item A (vraie) : au centre, $R=0$ donc $v=\omega R=0$.
  \item C (vraie) : rotation d'un solide $\Rightarrow$ même $\omega$ partout.
  \item D (vraie) : $v=\omega R$ croît avec $R$, donc $v_C>v_B>v_A$.
\end{itemize}
\end{corrige}

\begin{corrige}[la fronde : capstone]
\begin{enumerate}[label=\alph*)]
  \item \num{2} tours/s $\Rightarrow T=\qty{0,5}{\second}$ ; $\omega=\dfrac{2\pi}{T}=\dfrac{2\times\num{3,14}}{\num{0,5}}\approx\qty{12,56}{\radian\per\second}$.
  \item $v=\omega R=\num{12,56}\times\num{0,4}\approx\qty{5,02}{\metre\per\second}$.
  \item $a_c=\dfrac{v^2}{R}=\dfrac{(\num{5,02})^2}{\num{0,4}}\approx\qty{63}{\metre\per\second\squared}$ ;
        $F_c=m\,a_c=\num{0,3}\times63\approx\qty{19}{\newton}$.
  \item La pierre part en \textbf{ligne droite}, tangente au cercle au point de rupture, à la vitesse
        qu'elle avait : c'est le \emph{principe d'inertie} ($1^{\text{re}}$ loi). Il n'y a pas de force qui la projette « vers l'extérieur ».
\end{enumerate}
\end{corrige}

\begin{corrige}[le virage d'une voiture]
$F_c=\dfrac{mv^2}{R}$.
\begin{enumerate}[label=\alph*)]
  \item $F_c\propto v^2$ : en doublant $v$, la force est multipliée par $2^2=\mathbf{4}$.
  \item $F_c\propto \dfrac{1}{R}$ : en divisant $R$ par $2$, la force est multipliée par $\mathbf{2}$.
\end{enumerate}
\end{corrige}
\end{document}
